\begin{textsample}{POS Dim 6 – vem\_ed\_32 – Score 4.00 – vem\_ed\_32/t168.txt}  \label{ex:f6_pos_007}
IVEC 2025: pesquisa e inovação em Intercâmbios Virtuais A equipe de Projetos Colaborativos Internacionais (PCIs) da Divisão de Extensão e Pesquisa (Depes) da Coordenadoria de Ensino Superior de Graduação (CGESG) do Centro Paula Souza (CPS) participou da International Virtual Exchange Conference (IVEC) 2025, entre os dias 15 e 17 de outubro.
O evento híbrido é considerado um dos mais importantes do mundo na área de intercâmbios virtuais / COIL.
Voltada a professores, gestores, pesquisadores, estudantes e líderes educacionais, a conferência IVEC trouxe \textbf{programação} extensa de palestras, mesas-redondas, encontros temáticos, \textbf{workshops} e \textbf{apresentações} orais.
A instituição anfitriã desta edição foi a Hellenic Mediterranean University (Creta, Grécia). “Este encontro global oferece uma plataforma para a troca de conhecimentos e a compreensão de novas perspectivas pedagógicas.
Para o CPS, é uma oportunidade de destacar as inovações implementadas na gestão de seus Projetos Colaborativos Internacionais, consolidando a instituição como referência e inspirando novas colaborações”, comenta Osvaldo Succi Junior, coordenador da Área de Apoio à Internacionalização do Ensino Superior da Depes / CGESG, que participou de dois encontros temáticos virtuais: “COIL Connect: Building a sustainable research group” - com os professores Jon Rubin e Daisy Ruedas (COIL Connect, EUA) e Simone Hackett (The Hague University of Applied Sciences, Países Baixos) * A peer-driven approach in COIL Innovation: The ‘Magic Makers’ Experience for COIL Coordinators to Thrive in a Dynamic and Evolving Landscape” - com as professoras Gabriela Méndez Carrera, consultora independente, e Rosi León (DePaul University, EUA) As professoras Neusa Haruka Sezaki Gritti e Regiane Camargo, responsáveis, respectivamente, pelos PCIs em inglês e espanhol conduzidos pela Depes/ CGESG nas Fatecs, apresentaram o \textbf{workshop} online “911: Tecnologia ao resgate frente a linguistic tensions e intercultural challenges”, com os professores Alejandro Molina (Universidad Católica Silva Henríquez, Chile) e Guadalupe Vadillo (Universidad Nacional Autónoma de México - UNAM).

% matched lemmas: apresentação, programação, workshop
\end{textsample}
